\chapter{Aufbau und Gliederung der Arbeit}
\label{cha:Aufbau}

Der Aufbau, die Gliederung und der Umfang einer wissenschaftlichen Arbeit variiert ja nach Art und Thema der Arbeit. Das folgende Gliederungsschema hat sich allerdings bewährt und sollte im Wesentlichen eingehalten werden:
%
\begin{itemize}
	\item Deckblatt und Titelseite
%	\item	\emph{Sperrvermerk}
	\item Eidesstattliche Erklärung
	\item Aufgabenstellung
	\item Kurzfassung und Abstract
	\item Inhaltsverzeichnis
	\item Einleitung
	\item Grundlagen
	\item Hauptteil
	\item Zusammenfassung und Ausblick
	\item \textit{Abbildungs-, Tabellen-, Symbol- und Abkürzungsverzeichnis}
	\item Literaturverzeichnis
	\item \textit{Anhänge}
\end{itemize}
%
Die kursiv gesetzten Einträge sind dabei als optional zu betrachten. In den folgenden Abschnitten wird kurz auf die einzelnen Gliederungspunkte eingegangen.

Es ist grundsätzlich darauf zu achten, dass nicht mehr als drei Gliederungsebenen verwendet werden, also Kapitel, Unterkapitel und Unter-Unterkapitel. Wenn weitere Gliederungsebenen erforderlich werden, ist das grundsätzliche Konzept der Ausarbeitung zu überdenken und entsprechend abzuändern.

\section{Deckblatt und Titelseite}
\label{sec:Titelseite}

Das Deckblatt und die Titelseite sind die ersten beiden Seiten dieses Leitfadens bzw. dieser Vorlage. Sie müssen entsprechend gestaltet werden und alle darauf aufgeführten Informationen enthalten. Auf dem Deckblatt sind das das Logo des Instituts, das Logo der Leibniz Universität Hannover, der Name der Autorin bzw. des Autors, der Titel der Arbeit, die Inventar-Nummer der Arbeit (diese wird kurz vor Abgabe der Arbeit vom Institut vergeben und ist in die Box zu setzen) sowie die Anschrift des Instituts. Auf der Titelseite sind der Passus \textit{Ausarbeitung für das Modul Wissenschaftliches Arbeiten im Ingenieurwesen}, der Studiengang, die Matrikelnummer, das Abgabedatum sowie der Prüfer und die Betreuerin bzw. der Betreuer gemäß dieser Vorlage zu ergänzen.

Das Deckblatt ist außen auf den Umschlag der gebundenen, gedruckten Ausgabe aufzudrucken. Die Titelseite bildet die erste Papierseite, die in den Umschlag eingebunden wird. Anschauungsexemplare können auf Anfrage am Institut für Windenergiesysteme eingesehen werden.

%\section{Sperrvermerk}
%\label{sec:Sperr}

%Bei wissenschaftlichen Arbeiten, die in Zusammenarbeit mit einem Unternehmen erstellt werden, kann ein Sperrvermerk der Arbeit gefordert werden. Der Sperrvermerk schützt dabei die vertraulichen Daten des Unternehmens gegenüber der Öffentlichkeit (z.B. Wettbewerber des Unternehmens) und ist als erste Seite unmittelbar nach dem Deckblatt in die Arbeit einzubinden. Optional kann der Sperrvermerk direkt auf dem Deckblatt platziert werden.

%Der Inhalt des Sperrvermerkes ist mit dem Betreuer abzustimmen.

\section{Eidesstattliche Erklärung}
\label{sec:Eid}

Eine eidesstattliche Erklärung ist ein wichtiger Bestandteil einer wissenschaftlichen Arbeit und muss der Arbeit vorangestellt werden. In der eidesstattlichen Erklärung versichert die Verfasserin bzw. der Verfasser, dass die vorliegende Arbeit eigenständig und ohne die Hilfe Dritter angefertigt wurde. Des Weiteren wird erklärt, dass alle verwendeten Quellen als solche kenntlich gemacht und im Literaturverzeichnis aufgeführt wurden. Außerdem wird durch den Verfasser versichert, dass die Arbeit in dieser Form weder in Deutschland noch im Ausland bereits veröffentlicht oder für ein Studium vorgelegt oder eingebracht wurde. Insbesondere wird auch bestätigt, dass künstliche Intelligenz zur Erstellung von Inhalten entweder nicht eingesetzt oder entsprechende Passagen kenntlich gemacht wurden. Wie der Einsatz von KI-Tools nach der Handreichung der Leibniz Universität Hannover dokumentiert werden kann, zeigt Anhang~\ref{cha:KI-Dokumentation}.

Der Wortlaut der eidesstattlichen Erklärung muss dem auf Seite \pageref{Erklaerung} entsprechen. Die Erklärung ist mit Ort, Datum sowie der Unterschrift der Verfasserin bzw. des Verfassers abzuschließen.

\section{Aufgabenstellung}
\label{sec:Aufgabe}

Die Aufgabenstellung gibt Auskunft über das Thema der Arbeit, beinhaltet die zu bearbeitenden Teilaufgaben und dokumentiert das Aus- und das spätestmögliche Abgabedatum sowie die Namen des Prüfers und der Betreuerin bzw. des Betreuers. Die Aufgabenstellung wird vom Institut für Windenergiesysteme ausgegeben. Es soll im Original hinter der Eidesstattlichen Erklärung eingebunden werden. Die zwei Seiten für die Aufgabenstellung dieser Vorlage sind als Platzhalter zu verstehen und sind durch die tatsächlich ausgegebene Aufgabenstellung zu ersetzen. 

\section{Kurzfassung und Abstract}
\label{sec:Abstract_DE}

Auf maximal einer halben Seite ist eine Kurzfassung der Arbeit voranzustellen. Die Kurzfassung soll der Leserin bzw. dem Leser einen Überblick über den Inhalt des Dokuments geben. Daher sollte sie eine kurze Einleitung ins Thema, die methodischen Ansätze, die zur Bearbeitung des Themas herangezogen wurden sowie die wesentlichen Ergebnisse der Bearbeitung beinhalten. Eine englischsprachige Kurzfassung (Abstract) ist nach der deutschsprachigen Kurzfassung zu ergänzen. Beide Kurzfassungen sollen insgesamt den Umfang einer Seite nicht überschreiten.

\section{Inhaltsverzeichnis}
\label{sec:Inhalt}

Das Inhaltsverzeichnis umfasst alle im Dokument vorhandenen Kapitel, Unterkapitel und Unter-Unterkapitel. Auch weitere Verzeichnisse (Abbildungs-, Tabellen-, Symbol- und Abkürzungsverzeichnisse) und ggf. Anhänge müssen im Literaturverzeichnis enthalten sein. 

Die Überschriften im Text und im Inhaltsverzeichnis müssen exakt übereinstimmen. Überschriften einzelner Gliederungspunkte dürfen sich nicht mit dem Gesamtthema der Arbeit und nicht mit Überschriften von Unterpunkten decken. Es ist darauf zu achten, dass aus den Überschriften der zu erwartende Inhalt hervorgeht. Allgemeingültige Kapiteltitel wie z. B. \textit{Grundlagen} oder \textit{Hauptteil} sind zu vermeiden.

Neben den durchnummerierten Abschnittsbezeichnungen ist die Seitenzahl des Abschnittsbeginns aufzuführen.

\section{Einleitung}
\label{sec:Einleitung}

Die Einleitung bildet das erste nummerierte Kapitel der Ausarbeitung. Sie leitet in das Thema ein, motiviert die Arbeit und leitet die Problemstellung ab. Die Einleitung dient dazu, das Thema in den übergeordneten Kontext zu setzen und ggf. eine Abgrenzung des Themas vorzunehmen. Den Abschluss der Einleitung bildet eine Beschreibung des Aufbaus der Ausarbeitung.

\section{Technisch-physikalische Grundlagen}
\label{sec:Grundlagen}

Das zweite Kapitel einer wissenschaftlichen Arbeit widmet sich den technisch-physikalischen Grundlagen, die ein fachlich gebildeter Leser kennen muss, um alle folgenden Kapitel nachvollziehen zu können. Es ist darauf zu achten, diesem Kapitel einen inhaltsbezogenen Titel zu geben. Ein allgemeiner Titel wie \textit{Technisch-physikalische Grundlagen} ist nicht zulässig.

\section{Hauptteil}
\label{sec:Hauptteil}

Im Hauptteil einer wissenschaftlichen Arbeit wird die Gesamtheit der bearbeiteten Problemstellung dargestellt. Die methodischen Ansätze sowie der Lösungsweg sind ebenso darzustellen wie die Ergebnisse der Untersuchungen. 

Generell ist auf eine klare Darstellung des Gedankengangs, eine der Nachvollziehbarkeit der Ausarbeitung dienliche Gliederung und folgerichtige Anordnung des Inhalts sowie auf eine genaue, verständliche, kurze und prägnante Formulierung zu achten.

Die Gliederung des Inhalts und damit die Reihenfolge der Darstellung müssen sich nach sachlichen Gesichtspunkten und nicht nach der zeitlichen Abfolge der Bearbeitung richten. Der Hauptteil kann mehrere thematisch getrennte Kapitel umfassen. Es ist darauf zu achten, alle Kapitel mit inhaltsbezogenen Titeln zu versehen. Ein allgemeingültiger Titel wie \textit{Hauptteil} ist nicht zulässig.

Durch die Beschreibung der durchgeführten Untersuchungen müssen die erzielten Ergebnisse reproduzierbar sein. Werden Ergebnisse oder Lösungen von anderen Quellen zitiert, so sind diese nur soweit zu erläutern, wie es für das Verständnis notwendig ist. Eine Herleitung kann hingegen entfallen. Allerdings sollten Herleitungen zu tatsächlich selbst verwendeten Methoden in jedem Fall eingebunden werden.

\section{Zusammenfassung und Ausblick}
\label{sec:Zusammenfassung}

Eine wissenschaftliche Arbeit wird durch eine Zusammenfassung abgeschlossen. Die Zusammenfassung hat klare Aussagen über die verwendeten Verfahren, die erhaltenen Ergebnisse und ihre Bedeutung, die Grenzen der Gültigkeit und die Fortschritte gegenüber dem Kenntnisstand zu Beginn der Arbeit sowie ggf. die Anwendbarkeit der Ergebnisse zu enthalten. 

Des Weiteren ist in der Zusammenfassung in geeigneter Weise auf den Eigenanteil der Arbeit hinzuweisen und eventuell ein Ausblick auf mögliche vertiefende Forschungsarbeiten oder weiterführende Untersuchungen der Thematik zu geben.

Eine Leserin bzw. ein Leser, der auf der Suche nach Informationen zu einem bestimmten Thema ist, sollte nach Durchsicht der Kurzfassung, der Zusammenfassung und der Einleitung einen guten Überblick über die gesamte Arbeit bekommen haben. Daher werden in der Zusammenfassung in der Regel Aussagen wiederholt, die im Hauptteil der Arbeit bereits ausführlich beschrieben worden sind.

\section{Abbildungs, Tabellen-, Symbol- und Abkürzungsverzeichnis}
\label{sec:Abb-Tab-Symb-Abk-v}

Das Abbildungsverzeichnis hat den Zweck, alle Abbildungen bzw. deren Bildunterschriften mit Angabe der Seitenzahlen aufzulisten. Das Abbildungsverzeichnis ist optional, muss allerdings eingefügt werden, wenn mehr als drei Abbildungen im Dokument enthalten sind. Das Tabellenverzeichnis ist analog zum Abbildungsverzeichnis einzubinden, wenn mehr als drei Tabellen im Dokument vorhanden sind.

Das Tabellenverzeichnis hat den gleichen Zweck wie das Abbildungsverzeichnis und ist ebenfalls optional.

Das Symbolverzeichnis dient der Auflistung aller verwendeter Symbole. Dies ist vor allem für Arbeiten mit vielen Formeln hilfreich, da Formelzeichen als Symbole aufgefasst werden. Die Aufführung erfolgt tabellarisch und sollte Angaben zu den verwendeten Symbolen, ihrer Bedeutung und Einheit enthalten. Das Symbolverzeichnis ist alphabetisch zu sortieren. Griechische und römische Symbole sind in getrennten Blöcken aufzulisten.

Das Abkürzungsverzeichnis wird einer wissenschaftlichen Arbeit dann beigefügt, wenn im Fließtext viele nicht geläufige Abkürzungen verwendet werden. In tabellarischer Form werden sowohl die Abkürzungen als auch ihre Bedeutungen aufgelistet. Hierbei sind nur nicht im "`Duden"' verzeichnete Abkürzungen anzugeben. 

\section{Literaturverzeichnis}
\label{sec:Literatur}

Im Literaturverzeichnis werden alle zitierten Quellen aufgeführt. Dies hat in alphabetischer Reihenfolge des Nachnamens der Erstautorinnen bzw. Erstautoren zu erfolgen. Die Nummern der Quellen haben mit den im Text verwendeten Hinweisnummern auf die entsprechende Literatur oder Quelle übereinzustimmen.

Weitere Informationen bezüglich der Zitierweise werden in Abschnitt \ref{sec:Zitier} gegeben. Layoutbeispiele für verschiedene Quellenarten finden sich im Literaturverzeichnis am Ende des Dokuments. Der Stil des Literaturverzeichnisses sollte in jedem Fall übernommen werden.

\section{Anhänge}
\label{sec:Anhänge}

Anhänge in einer wissenschaftlichen Arbeit dienen der Darstellung von Inhalten wie Abbildungen, Tabellen, Herleitungen, Programmquelltexten, etc., die im Text keinen Platz finden oder dort unnötig sind und dennoch für die Nachvollziehbarkeit und die Reproduzierbarkeit der Arbeit von Bedeutung sind. Anhänge werden im Inhaltsverzeichnis aufgeführt oder durch ein separates Verzeichnis aufgelistet. Im Fließtext sollte auf die Bezug genommenen Anhänge verwiesen werden.

